\documentclass[11pt]{article}
\usepackage[T1]{fontenc}
\usepackage[utf8]{inputenc}
\usepackage{lmodern,microtype}
\usepackage[margin=1.08in]{geometry}
\usepackage{amsmath,amssymb,amsthm,mathtools}
\usepackage{xcolor,hyperref}
\hypersetup{colorlinks=true,linkcolor=blue!55!black,citecolor=blue!55!black,urlcolor=blue!60!black}
\setlength{\parindent}{0pt}
\setlength{\parskip}{0.48em}
\setlength{\emergencystretch}{2em}
\newtheorem{theorem}{Theorem}
\theoremstyle{remark}
\newtheorem{remark}[theorem]{Remark}

\title{A Counterexample at the Omitted Parameter in an Affine-Maximal Bernstein Statement}
\author{[Author Name]}
\date{}

\begin{document}
\maketitle

\begin{abstract}
\textbf{Relation to the problem list.}
This manuscript concerns Question~2 (L2), ``Affine-maximal Bernstein
threshold,'' but it is \emph{not} a solution of the intended Sun--Xu
conjecture.  The question in the list accidentally omits the source's standing
assumption $a\ne0$.  The counterexample below uses only the resulting
degenerate case $a=0$; the intended nonzero-parameter problem remains open.

Let $f$ be a smooth strictly convex function on a convex domain in
$\mathbb R^2$, let $D=\det D^2f$, and consider the affine-maximal-type
equation $f^{ij}(D^a)_{ij}=0$.  A proposed Bernstein statement asserts
quadratic rigidity for complete Calabi metrics when
$a\notin[-2/3,-1/3]$.  We show that the restriction $a\ne0$ is essential:
at $a=0$ the equation is vacuous, and the explicit nonquadratic potential
$f(x,y)=\frac12(x^2+y^2)+x^4$ is strictly convex and has a complete Calabi
metric.  Thus any formulation that includes $a=0$ is false.
\end{abstract}

\section{Introduction}

For a smooth strictly convex function $f\colon\Omega\to\mathbb R$, write
\[
 D=\det D^2f>0,
 \qquad
 L_f u=f^{ij}u_{ij},
\]
where $(f^{ij})=(D^2f)^{-1}$.  The affine-maximal-type equation studied by
Sun and Xu \cite{SunXu} is
\begin{equation}\label{eq:pde}
 L_f(D^a)=f^{ij}(D^a)_{ij}=0.
\end{equation}
The associated Calabi metric is the Hessian metric
\begin{equation}\label{eq:calabi}
 G_f=f_{ij}\,dx^i dx^j.
\end{equation}
Their Bernstein conjecture concerns nonzero parameters outside the interval
$[-2/3,-1/3]$.  If the standing condition $a\ne0$ is omitted, however,
the assertion degenerates immediately.  The following theorem records a
complete counterexample.

\begin{theorem}\label{thm:main}
There exists a smooth strictly convex nonquadratic function
$f\colon\mathbb R^2\to\mathbb R$ for which the Calabi metric $G_f$ is
complete and \eqref{eq:pde} holds with $a=0$.  Consequently, the Bernstein
statement allowing every $a\notin[-2/3,-1/3]$ is false.
\end{theorem}

\section{The counterexample}

\begin{proof}[Proof of Theorem~\ref{thm:main}]
Set
\begin{equation}\label{eq:f}
 f(x,y)=\frac12(x^2+y^2)+x^4.
\end{equation}
Then
\[
 D^2f=
 \begin{pmatrix}
  1+12x^2&0\\
  0&1
 \end{pmatrix},
 \qquad
 D=1+12x^2.
\]
Thus $D^2f$ is positive definite at every point, so $f$ is smooth and
strictly convex.  It is plainly nonquadratic.

At $a=0$ one has $D^a\equiv1$.  Therefore
\[
 L_f(D^0)=L_f(1)=f^{ij}(1)_{ij}=0,
\]
and \eqref{eq:pde} is satisfied identically.

It remains to verify completeness.  From $D^2f\ge I$ as quadratic forms,
the length of every piecewise $C^1$ curve $\gamma$ satisfies
\[
 \operatorname{Length}_{G_f}(\gamma)
 =\int\sqrt{D^2f(\dot\gamma,\dot\gamma)}\,dt
 \ge \int |\dot\gamma|\,dt
 =\operatorname{Length}_{\mathrm{Euc}}(\gamma).
\]
Hence the identity map
$(\mathbb R^2,G_f)\to(\mathbb R^2,g_{\mathrm{Euc}})$ is distance
nonincreasing.  Every $G_f$-Cauchy sequence is Euclidean Cauchy and converges
to a point of $\mathbb R^2$; on a Euclidean compact neighborhood of that
point, $G_f$ is uniformly equivalent to the Euclidean metric, so convergence
also holds in $G_f$.  Thus $(\mathbb R^2,G_f)$ is a complete metric space,
and therefore geodesically complete by Hopf--Rinow.  This proves the theorem.
\end{proof}

\begin{remark}
The construction does not address the nondegenerate parameter range
$a\ne0$.  Its point is that the zero parameter cannot be added to the
statement: at $a=0$, every smooth strictly convex potential satisfies the
differential equation, and completeness alone cannot force a quadratic.
\end{remark}

\begin{thebibliography}{9}

\bibitem{SunXu}
Y.~Sun and R.~Xu,
\emph{A class of new complete affine maximal type hypersurfaces},
Differential Geom. Appl. \textbf{101} (2025), 102308;
\href{https://arxiv.org/abs/2501.08525}{arXiv:2501.08525}.

\bibitem{Shima}
H.~Shima,
\emph{The Geometry of Hessian Structures},
World Scientific, Singapore, 2007.

\end{thebibliography}
\end{document}
